\section{年轻研究者如何进入架构研究}

前面已经把 Attention 的技术路线和硬件约束讲完，本章把讨论落回研究方法。架构研究不是单纯读论文，也不是只会写公式；它需要接近真实模型、真实算力和真实硬件约束。对于年轻研究者来说，最重要的问题是怎样进入真正能验证 scale 的环境。

节目结尾给年轻研究者的建议很直接：找一个 lab 实习，跟上 frontier，因为做架构必须有算力。这个建议听起来朴素，但背后有很现实的原因：架构研究很难只靠小规模实验判断，小模型上 work 的想法到大模型未必 work。

\begin{figure}[H]
\centering
\includegraphics[width=0.96\textwidth]{young-researcher-advice.png}
\caption{进入架构研究的路径：考古、算力、frontier lab 和硬件意识缺一不可。自制概念图，依据 01:42:23--01:43:26 对谈内容整理。}
\end{figure}

\begin{knowledgebox}{读图：架构研究需要靠近真实约束}
读旧论文建立地图，找到有用问题，获得足够算力验证 scale，进入 frontier lab 接近真实训练环境，再用开源实现让技术流传。这是一条从理解到验证再到影响的路径。
\end{knowledgebox}

\subsection{为什么算力是入场券}

架构研究的许多问题只有到一定规模才暴露。表达力、训练稳定、KV cache、decoding、kernel、batch size、硬件吞吐，这些都不是单机小实验能完全验证的。没有算力，就很难判断一个架构是未来路线，还是小规模幻觉。

\subsection{本章小结}

进入架构研究，既要有算法品味，也要有工程和硬件意识。年轻研究者不能只读最新论文，也不能只做玩具实验；最重要的是靠近真实 frontier 约束。

